\documentclass[12pt,a4paper]{article}
\usepackage[utf8]{inputenc}
\usepackage[T1]{fontenc}
\usepackage[french]{babel}
\usepackage{geometry}
\geometry{margin=2.5cm}
\usepackage{xcolor}
\usepackage{graphicx}
\usepackage{tikz}
\usetikzlibrary{shapes.geometric, arrows.meta, positioning, shadows, calc}
\usepackage{tabularx}
\usepackage{booktabs}
\usepackage{colortbl}
\usepackage{hyperref}
\usepackage{enumitem}
\usepackage{fancyhdr}
\usepackage{titlesec}
\usepackage{listings}

% Définition des couleurs
\definecolor{primary}{HTML}{2B6CB0} % Bleu professionnel
\definecolor{secondary}{HTML}{48BB78} % Vert
\definecolor{accent}{HTML}{ED8936} % Orange
\definecolor{bglight}{HTML}{F7FAFC} % Gris très clair pour les fonds
\definecolor{codegray}{HTML}{2D3748}

\hypersetup{
    colorlinks=true,
    linkcolor=primary,
    urlcolor=accent
}

% Configuration des titres
\titleformat{\section}{\Large\bfseries\color{primary}}{\thesection.}{0.5em}{}
\titleformat{\subsection}{\large\bfseries\color{secondary}}{\thesubsection.}{0.5em}{}

% Configuration de l'en-tête et pied de page
\pagestyle{fancy}
\fancyhf{}
\fancyhead[L]{\small Résumé Explicatif des Datasets}
\fancyhead[R]{\small Nexora}
\fancyfoot[C]{\thepage}

\begin{document}

\begin{titlepage}
    \centering
    \vspace*{2cm}
    
    {\Huge \bfseries \color{primary} Analyse et Résumé des Datasets}\\[1cm]
    {\Large \textit{Dossier \texttt{data} -- Projet Nexora}}\\[2cm]
    
    \begin{tikzpicture}
        \node[circle, fill=primary!10, draw=primary, line width=2pt, minimum size=4cm] (A) {\Large \textbf{DATA}};
        \node[circle, fill=secondary!10, draw=secondary, line width=2pt, minimum size=2.5cm] at (4,2) (B) {JSONL};
        \node[circle, fill=accent!10, draw=accent, line width=2pt, minimum size=2.5cm] at (4,-2) (C) {Python};
        \draw[thick, primary, dashed] (A) -- (B);
        \draw[thick, primary, dashed] (A) -- (C);
    \end{tikzpicture}
    
    \vspace{2cm}
    
    Ce document présente une analyse détaillée, visuelle et simplifiée de tous les fichiers présents dans le répertoire de données. Il illustre la structure, le volume et les relations entre les différentes entités métier générées pour la plateforme.
    
    \vfill
    {\large \textbf{Date :} \today}
\end{titlepage}

\tableofcontents
\newpage

%====================================================================
\section{Aperçu Global des Fichiers}
%====================================================================

Le dossier \texttt{data/} contient plusieurs fichiers de données au format \textbf{JSONL} (un objet JSON par ligne) ainsi qu'un script Python responsable de leur génération. L'utilisation du format JSONL est idéale pour le Big Data, car elle permet de lire les données ligne par ligne sans saturer la mémoire (streaming).

\subsection{Liste et Statistiques des Fichiers}

Le tableau ci-dessous dresse l'inventaire exact de ces fichiers :

\begin{table}[h!]
    \centering
    \renewcommand{\arraystretch}{1.5}
    \begin{tabularx}{\textwidth}{>{\columncolor{bglight}}l c l X}
        \toprule
        \rowcolor{primary}
        \textcolor{white}{\textbf{Fichier}} & \textcolor{white}{\textbf{Volume (Records)}} & \textcolor{white}{\textbf{Taille}} & \textcolor{white}{\textbf{Description simplifiée}} \\
        \midrule
        \texttt{companies.jsonl} & 50 & $\sim$33 Ko & Entreprises réelles et fictives avec leurs caractéristiques. \\
        \texttt{documents.jsonl} & 500 & $\sim$850 Ko & Base de connaissances (articles, guides, tutoriels). \\
        \texttt{employees.jsonl} & 500 & $\sim$1 Mo & Profils des experts informatiques et leurs compétences. \\
        \texttt{projects.jsonl} & 120 & $\sim$120 Ko & Projets de développement et membres de l'équipe. \\
        \texttt{skills.jsonl} & 82 & $\sim$8 Ko & Référentiel des compétences techniques requises. \\
        \texttt{technologies.jsonl}& $\sim$20 000 & $\sim$3 Mo & Extractions Wikipedia/DBpedia de concepts technologiques. \\
        \texttt{generate\_data.py} & N/A & $\sim$20 Ko & Script métier générant les données de manière réaliste. \\
        \bottomrule
    \end{tabularx}
    \caption{Résumé des fichiers contenus dans le dossier data.}
\end{table}

%====================================================================
\section{Modèle de Données et Relations}
%====================================================================

Afin de bien comprendre comment ces fichiers s'imbriquent pour former un graphe de connaissances, voici une illustration des relations logiques entre eux.

\vspace{1cm}
\begin{center}
\begin{tikzpicture}[
    node distance=2.5cm and 3cm,
    entity/.style={rectangle, draw=primary, fill=primary!5, ultra thick, minimum width=3.5cm, minimum height=1.5cm, text centered, rounded corners, font=\bfseries},
    relation/.style={->, >=stealth, very thick, color=gray!80!black, font=\footnotesize\itshape}
]

% Noeuds
\node[entity, fill=blue!10] (employee) {Employees};
\node[entity, fill=green!10, right=of employee] (skill) {Skills};
\node[entity, fill=orange!10, below=of employee] (project) {Projects};
\node[entity, fill=purple!10, left=of employee] (document) {Documents};
\node[entity, fill=teal!10, right=of project] (company) {Companies};

% Liens
\draw[relation] (employee) -- node[above] {possède (niveau)} (skill);
\draw[relation] (employee) -- node[right] {travaille sur} (project);
\draw[relation] (employee) -- node[above] {est auteur de} (document);
\draw[relation] (project) -- node[right] {requiert} (skill);
\draw[relation] (project) -- node[below] {appartient à} (company);
\draw[relation] (company) -- node[left] {utilise le stack} (skill);

\end{tikzpicture}
\end{center}
\vspace{0.5cm}

\textbf{Lecture du graphe :} 
Un \textbf{Employé} possède des \textbf{Compétences}. Il travaille sur des \textbf{Projets} (qui requièrent ces mêmes compétences) et rédige des \textbf{Documents} techniques. Les Projets sont associés aux \textbf{Entreprises}.

\newpage
%====================================================================
\section{Analyse Détaillée de Chaque Dataset}
%====================================================================

\subsection{1. Les Profils Experts (\texttt{employees.jsonl})}
Ce fichier est le cœur du système. Il contient 500 profils d'experts générés avec une grande cohérence temporelle et structurelle.

\begin{itemize}[label=\textcolor{primary}{\textbullet}]
    \item \textbf{Données clés :} ID, Nom, Email, Département, Rôle, Localisation, Date d'embauche, Années d'expérience.
    \item \textbf{Particularité :} Chaque profil contient une liste imbriquée de ses compétences (\texttt{skills}) avec un niveau de maîtrise (de 1 à 5). Le nombre de compétences est proportionnel aux \textit{années d'expérience}.
\end{itemize}

\subsection{2. La Bibliothèque Technique (\texttt{documents.jsonl})}
Une collection de 500 documents rédigés par les experts pour alimenter la plateforme en connaissances.

\begin{itemize}[label=\textcolor{primary}{\textbullet}]
    \item \textbf{Données clés :} Titre, Type (Whitepaper, Tutorial, Guide...), Sujet (DevOps, IA, Cloud...), Résumé (Abstract), Contenu textuel, Auteur (lié à un employé), Vues, Note (Rating).
    \item \textbf{Utilité :} Permet de tester les algorithmes de recherche en texte intégral (Full-Text Search) et le NLP (Traitement du Langage Naturel).
\end{itemize}

\subsection{3. Le Portefeuille Projets (\texttt{projects.jsonl})}
Il liste 120 projets d'envergure nécessitant des collaborations en équipe.

\begin{itemize}[label=\textcolor{primary}{\textbullet}]
    \item \textbf{Données clés :} Nom, Domaine (ex. HealthTech, Fintech), Budget, Priorité, Statut (Actif, Terminé...).
    \item \textbf{Relations :} Possède un tableau d'identifiants d'employés (\texttt{team\_members}) et une liste de technologies requises.
\end{itemize}

\begin{center}
\begin{tikzpicture}
    % Petit graphique illustratif pour project
    \fill[primary!20] (0,0) rectangle (6,3);
    \draw[thick, primary] (0,0) rectangle (6,3);
    \node[font=\bfseries] at (3,2.5) {Anatomie d'un Projet};
    \draw[thick, dashed] (0.5, 2.2) -- (5.5, 2.2);
    \node[anchor=west] at (0.5, 1.7) {1. Budget \& Domaine métiers};
    \node[anchor=west] at (0.5, 1.1) {2. Equipe (IDs Employees)};
    \node[anchor=west] at (0.5, 0.5) {3. Technologies Requises};
\end{tikzpicture}
\end{center}

\subsection{4. Le Répertoire des Entreprises (\texttt{companies.jsonl})}
Un fichier de 50 enregistrements détaillant des entreprises partenaires ou clientes, allant des startups aux multinationales.

\begin{itemize}[label=\textcolor{primary}{\textbullet}]
    \item \textbf{Données clés :} Nom, Industrie, Taille, Nombre d'employés, Siège social, Stack technologique, Avantages sociaux.
    \item \textbf{Utilité :} Permet d'offrir une dimension "B2B" (Business to Business) à la plateforme, en permettant des recommandations de profils vers des entreprises spécifiques.
\end{itemize}

\subsection{5. Le Lexique des Compétences (\texttt{skills.jsonl})}
Il s'agit du dictionnaire standard de la plateforme, avec 82 compétences.

\begin{itemize}[label=\textcolor{primary}{\textbullet}]
    \item \textbf{Catégories :} Langages de programmation, Frameworks, Cloud, Data \& AI, DevOps, Sécurité.
    \item \textbf{Score de Demande :} Chaque compétence est dotée d'un score (\texttt{demand}) qui reflète son utilité sur le marché du travail actuel.
\end{itemize}

\subsection{6. Le Grand Dictionnaire \texttt{technologies.jsonl}}
C'est le plus gros fichier contenant plus de 20 000 objets. C'est un export issu du Web sémantique (DBpedia/Wikipedia) qui sert à élargir le vocabulaire du moteur de recommandation (jeux vidéo associés, logiciels historiques, architectures matérielles).

%====================================================================
\section{La Génération Intelligente (\texttt{generate\_data.py})}
%====================================================================

La grande force de ce dossier repose sur le script Python. Il ne s'agit pas d'un générateur aléatoire classique, mais d'un \textbf{générateur algorithmique corrélé}.

\begin{center}
\begin{table}[h!]
    \centering
    \renewcommand{\arraystretch}{1.3}
    \begin{tabularx}{0.8\textwidth}{|X|X|}
    \hline
    \rowcolor{secondary!20}
    \textbf{Technique utilisée} & \textbf{Résultat obtenu} \\
    \hline
    \textbf{Seed Déterministe} & Chaque exécution donne \textit{exactement} les mêmes données (pour la reproductibilité). \\
    \hline
    \textbf{Affinité de Département} & Un employé en "Data Science" recevra majoritairement des compétences en "Python" et "IA", pas en "React". \\
    \hline
    \textbf{Corrélation d'Expérience} & Plus l'employé est ancien, plus sa liste de compétences est longue et son niveau élevé. \\
    \hline
    \end{tabularx}
\end{table}
\end{center}

\section{Conclusion}
Le dossier \texttt{data} est complet, professionnel et parfaitement calibré pour une application Big Data et Machine Learning. L'utilisation du format léger JSONL rend l'injection vers Spark, Kafka ou Neo4j extrêmement performante. La richesse des données (tableaux imbriqués, dates réalistes, textes générés) permet de réaliser toutes les opérations de la Business Intelligence moderne (recherche sémantique, classification, recommandations en graphe).

\end{document}
